% paper/sections/11e_interface_crossing.tex
%
% §11.5  不連続界面通過時の整合性
% ═══════════════════════════════════════════════════════════════════════════════
\subsection{不連続界面通過時の整合性}
\label{sec:interface_crossing}
% ═══════════════════════════════════════════════════════════════════════════════

§\ref{sec:force_balance}--§\ref{sec:advection_pressure_coupling} では
等密度・単相の条件下で NS ソルバの物理的整合性を確認した．
二相流への拡張では，密度比 $\rho_l/\rho_g \gg 1$ の界面を跨ぐ
圧力勾配の整合性が追加の検証課題となる．

\noindent\textbf{目的：}
密度比 $\rho_l/\rho_g = 1000$ の二相場において，
界面を跨ぐ速度連続性と圧力勾配の滑らかさを定量的に検証する．

\noindent\textbf{評価対象：}
GFM 不連続面処理（§\ref{sec:gfm}）による速度補正項 $\bnabla p / \rho$ の
界面両側での整合性．

\noindent\textbf{前提となる検証済み結果：}
\begin{itemize}[noitemsep]
  \item GFM Laplace 圧力跳び誤差 $1.8\%$（$N=128$）：§~\ref{sec:verify_gfm}（表~\ref{tab:gfm_laplace}）
  \item CCD-PPE 変密度 Poisson 格子収束 $\Ord{h^7}$：§~\ref{sec:verify_ppe}
\end{itemize}

\noindent\textbf{結論：}
密度比 $10^3$ における界面通過時の速度連続性の定量的検証は，
GFM を統合した全二相 NS ソルバの完成後に実施する（§\ref{sec:future_work}）．
第~\ref{sec:validation}章の静止液滴ベンチマーク（§\ref{sec:val_static_drop}）が
この検証を包含する初の統合テストとなる．
